\documentclass[11pt]{article}
\usepackage[utf8]{inputenc}
\usepackage{amsmath,amssymb,amsthm}
\usepackage{booktabs}
\usepackage[margin=2.5cm]{geometry}
\usepackage{hyperref}

\newtheorem{theorem}{Theorem}
\newtheorem{proposition}[theorem]{Proposition}
\newtheorem{corollary}[theorem]{Corollary}
\newtheorem{lemma}[theorem]{Lemma}
\newtheorem{conjecture}[theorem]{Conjecture}
\theoremstyle{definition}
\newtheorem{definition}[theorem]{Definition}
\newtheorem{observation}[theorem]{Observation}
\theoremstyle{remark}
\newtheorem{remark}[theorem]{Remark}

\newcommand{\D}{\mathrm{D}}
\newcommand{\qbin}[2]{\left[\begin{smallmatrix}#1\\#2\end{smallmatrix}\right]}
\newcommand{\gl}{\mathfrak{gl}}
\newcommand{\Sym}{\mathrm{Sym}}

\title{The defect of the differential expansion is a supergroup genus bound}
\author{\emph{(draft; authors to be added)}}
\date{}

\begin{document}
\maketitle

\begin{abstract}
The differential expansion of the colored HOMFLY polynomial $H_{[r]}(A,q)$ of a knot in symmetric
representations was conjectured to factorise with a defect $\delta=\deg\Delta-1$ determined by the Alexander
polynomial.  Using colored HOMFLY polynomials of all $2893$ knots with at most $12$ crossings and all
representations with at most six boxes, we find $19$ knots for which this fails, some of them with fractional
defect.  We describe the defect by the tower $e_m$ of $\lambda$-degrees of the specialisations
$P_m(q,\lambda)=H_{[r]}(q^{-m},q)|_{q^r=\lambda}$ and show that $P_m$ is the Links--Gould invariant
$LG^{m+1,1}$ of $U_q\gl(m+1|1)$, with continuous weight $\alpha=r$.  The genus bounds of Kohli--Tahar and of
L\'opez Neumann--Garoufalidis then give $(m+1)\deg\Delta\le e_m\le(m+1)g$ for every knot, so the conjecture holds
whenever the Alexander polynomial detects the genus (alternating and fibered knots), and the defect lies in
$[\deg\Delta-1,\,g-1]$ in general.  For arbitrary Young diagrams $R$ we show experimentally that $H_R$ on the line
$A=q^{-m}$ is the invariant of the covariant $U_q\gl(k|m+k)$-module $V_R$; this explains the universal reductions of
$H_R$ on these lines as atypicality and Berezinian twists, predicts new ones, and identifies the hook families
$[r,1^n]$ with the $\mathfrak{sl}(2|1)$-modules $V(n,a)$.  On every anomalous knot we could test, the hook family
$[r,1]$ attains the genus bound where the symmetric family does not: the anomalous defect is a property of the
Links--Gould member, not of the knot.
\end{abstract}

\tableofcontents

%======================================================================
\section{Introduction}

For a knot $K$ and a Young diagram $R$ let $H_R(A,q)$ be the reduced colored HOMFLY polynomial.  For symmetric
$R=[r]$ the \emph{differential expansion}
\begin{equation}\label{DE}
  H_{[r]}=1+\sum_{k=1}^{r}\qbin{r}{k}\,\D_{r+k-1}\cdots\D_{r}\;G_k(A,q),\qquad \D_n=Aq^n-A^{-1}q^{-n},
\end{equation}
has $r$-independent coefficients $G_k$.  In \cite{KM-defect} it was observed that the $G_k$ contain further
factors $\D_{-1}\D_0\cdots\D_{m_k-2}$ with $m_k=\lfloor\frac{k-1}{\delta+1}\rfloor+1$, and conjectured that the
\emph{defect} $\delta$ is $\deg\Delta-1$, where $\deg\Delta$ is half the span of the Alexander polynomial.

\paragraph{Results.}
\begin{enumerate}
\item (Section~\ref{sec:tower}, experimental.) On the whole table the factor structure of the $G_k$ is encoded in
integers $e_m$, $m\ge-1$, with $e_m+1$ the first $k$ for which $\D_m\mid G_k$; equivalently $e_m$ is a quarter of
the $\lambda$-span of $P_m(q,\lambda)=H_{[r]}(A=q^{-m},q)|_{q^r=\lambda}$.  The tower is superadditive, and its
slope $1+\delta=\lim e_m/(m+1)$ is the defect.  For $19$ knots $\delta\neq\deg\Delta-1$ (Table~\ref{tab:anom}),
including fractional slopes $3/2$.
\item (Theorem~\ref{thm:lg}.) $P_m$ is the Links--Gould invariant $LG^{m+1,1}$ at $\alpha=r$.  The proof is
super Schur--Weyl duality; we also verify it with an independent vertex model.
\item (Corollary~\ref{cor:bounds}.) Combined with \cite{KT,LNG}: $P_m(1,\lambda)=\Delta(\lambda^2)^{m+1}$ and
$(m+1)\deg\Delta\le e_m\le(m+1)g$ for every knot.  Hence $\delta=\deg\Delta-1$ for alternating and fibered knots, and
$\deg\Delta-1\le\delta\le g-1$ always.
\item (Section~\ref{sec:nonrect}, experimental.) For any $R$, $H_R(q^{-m},q)$ is the $(1,1)$-tangle invariant of
$V_R$ over $U_q\gl(k|m+k)$.  The universal reductions of $H_R$ on these lines are atypicality and
Berezinian twists, including new rules $(b_k)$.  Hook families $(r,\mu)$ give ``coloured Links--Gould''
polynomials $P^\mu_m$, and $P^{[1]}_1$ attains the genus bound on all anomalous knots tested
(Table~\ref{tab:hooks}).  This is a positive answer, on these examples, to Question~1.4 of \cite{LNG}, and proposes
a definition of the differential expansion for non-rectangular $R$.
\end{enumerate}

%======================================================================
\section{Conventions, data and methods}\label{sec:data}

$H_R$ is normalised by $A^{-1}H(L_+)-AH(L_-)=(q-q^{-1})H(L_0)$, $H_R(\text{unknot})=1$, topological framing;
$[n]=\frac{q^n-q^{-n}}{q-q^{-1}}$, $\qbin rk=\frac{[r]!}{[k]![r-k]!}$.  Then
$H_R(A,q)=H_{R^T}(A,-q^{-1})$, and $H_{[1]}(1,q)=\Delta(q^2)$.

\paragraph{Data.}  All $R$ with $|R|\le6$ for knots up to $11$ crossings, $|R|\le4$ (partly $\le6$) for $12$
crossings; symmetric $[r]$ up to $r=10$ for the arborescent knots $11n_{67},11n_{97},12n_{51},12n_{268},12n_{293},
12n_{457}$ and $r\le9$ for several knots with $\deg\Delta=2$.  The methods are the following.
\begin{itemize}
\item The Hecke algebra in seminormal form, for the fundamental representation.
\item Tangle calculus with inclusive and exclusive Racah matrices, for arborescent knots.
\item Cabling in the space of paths.
\end{itemize}
Each polynomial is reconstructed by dense interpolation modulo one or two $31$-bit primes.  It is then checked at
fresh points modulo an independent prime and at $q=1$.  Every anomalous case was cross-checked by cabling the
KnotInfo braid word.  Code and data are in the accompanying repository: \texttt{data/homfly}; the scripts named
below are in \texttt{scripts/}.

\paragraph{Independent vertex model.}  To test the supergroup statements we use a separate implementation that
does not use any of the methods above.
\begin{itemize}
\item The Perk--Schultz $R$-matrix of $U_q\gl(k|l)$ on $(\mathbb{C}^{k|l})^{\otimes N}$.
\item Projection to $V_R$ by Jucys--Murphy spectral projectors (for symmetric $R$, the joint $q$-eigenspace of
the Hecke generators).
\item The pivotal element and the partial trace.
\end{itemize}
It runs in floating point at random $q$ (\texttt{supergroup\_vertex\_check.py}, \texttt{nonrect\_supergroup.py}).
It reproduces $H_R(A=q^{l-k},q)$ in the mirror convention $q\mapsto q^{-1}$.

%======================================================================
\section{The tower $e_m$}\label{sec:tower}

\begin{observation}\label{obs:plus}
For all $2893$ knots and $k\le4$, and for $k\le10$ where data exist, the $G_k$ in (\ref{DE}) are Laurent
polynomials.
\end{observation}

\begin{observation}\label{obs:pers}
For all knots in the table:
\begin{itemize}
\item[(i)] if $\D_n\mid G_k$ then $\D_n\mid G_{k'}$ for $k'>k$;
\item[(ii)] if $\D_n\mid G_k$ then $\D_{n-1}\mid G_k$;
\item[(iii)] $\D_n^2\nmid G_k$.
\end{itemize}
\end{observation}

So the factor structure is encoded in $e_m+1:=\min\{k:\D_m\mid G_k\}$.  Equivalently, on the line $A=q^{-m}$ the
sum (\ref{DE}) stops at $k=e_m$, and
\[
  P_m(q,\lambda):=H_{[r]}(q^{-m},q)\big|_{q^r=\lambda}
  =1+\sum_{k=1}^{e_m}\qbin rk\prod_{i=0}^{k-1}\big(\lambda q^{i-m}-\lambda^{-1}q^{m-i}\big)\,G_k(q^{-m},q)\Big|_{q^r=\lambda}
\]
is a Laurent polynomial with $e_m=\frac14\mathrm{span}_\lambda P_m$.  Here $e_{-1}=0$ ($H_{[r]}(q,q)=1$) and
$P_0=\Delta(\lambda^2)$, so $e_0=\deg\Delta=:d$.

\begin{observation}\label{obs:superadd}
$f(n):=e_{n-1}$ is superadditive (tested on $2217$ knots with enough data).  By Fekete's lemma
$1+\delta:=\lim f(n)/n=\sup f(n)/n$ exists; we take it as the definition of the defect.
\end{observation}

\begin{table}[h]
\centering
\begin{tabular}{llcc l}
\toprule
knots & type & $d$ & $g$ & $e_{-1},e_0,e_1,e_2,\dots$\\
\midrule
$4_1$ (and all $d=g=1$) & regular & $1$ & $1$ & $0,1,2,3,4,5$\\
$6_2,\ 10_{129}$ & regular & $2$ & $2$ & $0,2,4,6,8$\\
$11n_{67},11n_{97},12n_{51},12n_{268}$ & I & $1$ & $2$ & $0,1,3,4,6,7,9$ \quad $(e_m=\lfloor\frac{3(m+1)}2\rfloor)$\\
$12n_{293},12n_{457}$ ($12n_{321}$, \dots) & II & $1$ & $2$ & $0,\mathbf 1,4,6,8$ \quad ($=2(m+1)$ for $m\ge1$)\\
$11n_{34},11n_{42},12n_{313},12n_{430}$ & $\Delta=1$ & $0$ & $3,2,2,2$ & $0,\mathbf 0,3,5$\\
$12n_{750},12n_{830}$ & & $2$ & $3$ & $0,2,6,9$\\
$11n_{73}$ & & $2$ & $3$ & $0,2,4,6,8$\\
$12n_{132}$ & & $2$ & $3$ & $0,2,4,\ge8$\\
\bottomrule
\end{tabular}
\caption{Towers $e_m$ of the anomalous knots (bold: the Alexander member lies below the slope).  The values
$e_1=6$, $e_2=9$ for $12n_{750},12n_{830}$ and $e_2=5$ for the $\Delta=1$ knots come from the vertex model
(Section~\ref{sec:nonrect}); the HOMFLY data alone only give lower bounds there.}\label{tab:anom}
\end{table}

\paragraph{Remarks on Table~\ref{tab:anom}.}
\begin{itemize}
\item All anomalous knots have $g>d$; the converse fails ($11n_{73}$).
\item Type II: only the Alexander member is degenerate, and for $m\ge1$ the tower saturates $e_m=(m+1)g$.
\item Type I: slope $3/2$, so $\delta=\frac12$.
\item $\Delta=1$ (formerly ``$\delta=-1$''): $e_1=3$.  The mutants $11n_{34}$ ($g=3$) and $11n_{42}$ ($g=2$) have
identical symmetric HOMFLY, so the symmetric tower cannot see $g$.
\item The slope is not determined by $(\tau,\sigma)$: for $g=2$, $d=1$, the Type I knots have $\tau=\sigma=0$,
but so does $12n_{23}$, which has $e_1=4$; and $12n_{63}$ ($\tau=-1,\sigma=2$, $g=3$) has $e_1=4<2g$.
\end{itemize}

%======================================================================
\section{$P_m$ is the Links--Gould invariant}\label{sec:lg}

Let $LG^{n,1}$ denote the invariant of $U_q\gl(n|1)$ on its $2^n$-dimensional typical modules $V_\alpha$
\cite{DKL,MW}.  For $n=2$ we use the variables $t_0^{1/2}=q^{-\alpha}$, $t_1^{1/2}=q^{\alpha+1}$ of \cite{KT}.

\begin{theorem}\label{thm:lg}
For every knot, $m\ge1$ and integer $r\ge m+1$, $H_{[r]}(q^{-m},q)$ is the scalar by which the $(1,1)$-tangle
of the knot acts on $\Sym^r\mathbb{C}^{1|m+1}$ (mirror convention).  This module is typical of dimension
$2^{m+1}$.  Hence $P_m=LG^{m+1,1}$ with $\lambda=q^\alpha$; in particular
\[
  P_1(q,\lambda)=LG^{2,1}\big(t_0=\lambda^{-2},\ t_1=q^2\lambda^2\big)\quad\text{(up to mirror image).}
\]
\end{theorem}

\begin{proof}
\emph{Step 1 (reduced HOMFLY as a Hecke scalar).}  Let $\beta\in B_n$ close to $K$ and let $\beta_R$ be its
$R$-cable, an element of the Hecke algebra $H_{n|R|}(q)$ after the substitution $\sigma_i\mapsto A\,T_i$ (framing
corrected).  Let $E:H_{N+1}\to H_N$ be Ocneanu's conditional expectation.  It is
$H_N$-bilinear, with $E(T_N)=1$ up to normalisation and $E(1)=\frac{A-A^{-1}}{q-q^{-1}}$.  Then
$E^{(n-1)|R|}(e_R\beta_Re_R)\in e_RH_{|R|}e_R=\mathbb{C}e_R$, for the primitive idempotent $e_R$, and the scalar is
$H_R(A,q)$.  The coefficients of $E$ are polynomial in $A^{\pm1}$ and $\frac{A-A^{-1}}{q-q^{-1}}$, so this identity
can be specialised to $A=q^{-m}$.  Note that $\dim_q[r]$ vanishes there for $r\ge2$, so the closed invariant is
$0$ and only the reduced one carries information.

\emph{Step 2 (super Schur--Weyl duality).}  On $V^{\otimes N}$, $V=\mathbb{C}^{1|m+1}$, the Perk--Schultz matrix gives a
representation of $H_N(q)$ commuting with $U_q\gl(1|m+1)$.  Under it, $e_{[r]}$ projects onto $\Sym^rV$.  The partial
quantum supertrace with the pivotal element is a Markov conditional expectation with the parameters of $E$ at
$A=q^{-m}$, since the quantum superdimension of $V$ is $\frac{q^{-m}-q^m}{q-q^{-1}}$.  A Markov expectation is
determined by these parameters, so it agrees with $E$ on the image.  Hence the $(1,1)$-tangle acts on $\Sym^rV$
by $H_{[r]}(q^{-m},q)$.

\emph{Step 3 (the module).}  $\Sym^rV$ has basis $e_0^{r-|S|}\prod_{i\in S}e_i$, $S\subseteq\{1,\dots,m+1\}$.  So
for $r\ge m+1$ its dimension is $2^{m+1}$, and its highest weight is $(r\,|\,0,\dots,0)$, which is atypical
exactly for $r\le m$.  The parity change $\gl(1|m+1)\cong\gl(m+1|1)$ mirrors the invariant.  Under it,
$\Sym^rV$ becomes the typical module $V_\alpha$ with $\alpha=r$, up to tensoring with a one-dimensional module, which
changes only the framing factor.

\emph{Step 4 (continuation).}  Both $P_m$ and $LG^{m+1,1}$ are Laurent polynomials in $\lambda=q^\alpha$ (and $q$)
that agree at $\lambda=q^r$ for all $r\ge m+1$, hence they are equal.
\end{proof}

\begin{remark}
Steps 1--2 are standard, but the details (normalisations, the framing twist in Step 3) should be written out in
full; we have not done so here.  They are confirmed numerically by the independent vertex model.  For
$\gl(1|2)$ with $r=2,3$ it matches $H_{[r]}(q^{-1},q)$ on $255$ knots: all knots up to $10$ crossings of braid
index $\le5$, together with the anomalous knots of braid index $\le5$.  It also matches for $r=4$ on
$3_1,4_1,5_1,5_2,7_1$, and for $\gl(1|3)$, $r=3$ (dimension $8$) on $62$ knots.  The relative error is below
$10^{-10}$.  We have not compared with an independent evaluation of $LG^{2,1}$; the trefoil gives
$1-t_0-t_1+t_0^2+2t_0t_1+t_1^2-t_0^2t_1-t_0t_1^2$.
\end{remark}

\begin{observation}\label{obs:lgchecks}
On all $1068$ knots whose data determine $P_1$ (completely for $302$ of them; \texttt{links\_gould\_check.py}):
\begin{itemize}
\item[(i)] $P_1(q,\lambda)=P_1(q,q/\lambda)$, the $t_0\leftrightarrow t_1$ symmetry;
\item[(ii)] $P_1(1,\lambda)=\Delta(\lambda^2)^2$;
\item[(iii)] $P_1(i,\lambda)=\pm\Delta(\lambda^4)$, matching $LG(t_0,-t_0^{-1})=\Delta(t_0^2)$;
\item[(iv)] $P_1(11n_{34})=P_1(11n_{42})$, as for $LG$ of the Conway and Kinoshita--Terasaka knots.
\end{itemize}
\end{observation}

\begin{corollary}\label{cor:bounds}
For every knot and $m\ge0$:
\[
 P_m(1,\lambda)=\Delta(\lambda^2)^{m+1},\qquad (m+1)\,d\le e_m\le(m+1)\,g,\qquad d-1\le\delta\le g-1 .
\]
If $\Delta$ detects the genus ($d=g$; e.g.\ alternating, in particular two-bridge, or fibered knots), then
$e_m=(m+1)g$ for all $m$ and $\delta=\deg\Delta-1$.
\end{corollary}
\begin{proof}
By Theorem~\ref{thm:lg}, $P_m$ is the invariant of a typical module of $\mathfrak{sl}(1|m+1)$, which has $m+1$
positive odd roots, with $t=q^{2\alpha}=\lambda^2$.  The specialisation $J(t,1)=\Delta(t)^{m+1}$ and the bounds
$(m+1)\deg_t\Delta\le\deg_tJ\le2(m+1)g$ are \cite[Thms.~1.1, 1.2, Cor.~1.3]{LNG}.  Since
$\deg_tJ=\frac12\mathrm{span}_\lambda P_m=2e_m$, this gives the bounds on $e_m$.  For $m=1$ it is \cite{KT}
($\mathrm{span}\,LG^{2,1}\le4g$, with equality for alternating knots).
\end{proof}

The data agree: on the $362$ two-bridge knots up to $12$ crossings, $e_m=(m+1)g$ wherever $P_m$ is complete
($124$ knots for $m=1$, $18$ for $m=2$, $18$ for $m=3$), and no incomplete case exceeds it.  The original conjecture
can fail only for $d<g$, and need not fail there ($11n_{73}$).

\begin{observation}[roots of unity]\label{obs:roots}
Let $q$ be a primitive $2k$-th root of unity, $2\le k\le m+1$, and write $m+1=a+kb$ with $b=\lfloor\frac{m+1}k\rfloor$.
Then $P_m(q,\lambda)=\pm\Delta(\lambda^2)^a\Delta(\lambda^{2k})^b$.  This was checked for $m=1$ on all $302$
knots with complete $P_1$, for $m=2$ on the $27$ with complete $P_2$, and for $m=3$ on $3_1$.  It fails for
$k=m+2$.
\end{observation}
This should follow from $H_{[r+k]}=H_{[r]}H_{[k]}$ at $q^{2k}=1$ and the duality of \cite{KM-roots}.  It forces
the coefficients of $\lambda^j$, $j>2(m+1)d$, to be divisible by $\Phi_1\prod_{k=1}^{m+1}\Phi_{2k}(q)$: e.g.\
$(q^2-1)(q^2+1)$ for $11n_{67}$, $m=1$.

%======================================================================
\section{Non-rectangular representations}\label{sec:nonrect}

For rectangular $R$, $H_R-1$ is divisible by $\D_{R_1}\D_{-\ell(R)}$; for non-rectangular $R$ by no $\D_n$.  On the
lines $A=q^{-n}$ there are instead universal reductions.
\begin{observation}[\texttt{line\_reductions.py}]\label{obs:red}
On $A=q^{-n}$, $n\ge0$ (for $n<0$ apply the rule to $R^T$):
\begin{itemize}
\item[(a)] if $n\ge R_1$, then $H_R=H_{R^*}$, with $R^*$ the transpose of the complement of $R^T$ in the
$n\times\ell(R)$ box;
\item[(b)] if $n<R_1$ and $R_2=n+1$, then $H_R=H_{(R_1+1,R_3,\dots)}$.
\end{itemize}
For $|R|\le6$, $|n|\le8$ and all knots, these are all pairwise coincidences, apart from (b$_2$) below and the
$\Delta=1$ knots at $A=1$.
\end{observation}

\subsection{The supergroup picture}
Let $V_R\subset(\mathbb{C}^{k|m+k})^{\otimes|R|}$ be the covariant $U_q\gl(k|m+k)$-module of $R$.  It is nonzero iff
$R_{k+1}\le m+k$.
\begin{observation}\label{obs:h1}
On $A=q^{-m}$, $H_R$ is the $(1,1)$-tangle invariant of $V_R$ for $\gl(k|m+k)$, for every $k$ with $V_R\neq0$.
With $V_{R^T}$ instead it fails.  This was checked by the vertex model on $4982$ pairs $(R,K)$:
\begin{itemize}
\item all $|R|\le6$ with $V_R\ne0$, up to the module sizes the code can handle;
\item $K$ ranging over all knots up to $9$ crossings and $14$ anomalous knots;
\item the algebras $\gl(1|1),\gl(1|2),\gl(1|3),\gl(2|2),\gl(2|3)$;
\item relative tolerance $10^{-7}$.
\end{itemize}
\end{observation}
The proof of Theorem~\ref{thm:lg} applies verbatim once Step 3 is replaced by the corresponding statement for
$V_R$.  For $k=1$, $V_R$ has highest weight $(R_1\,|\,\mu^T)$ with $\mu=(R_2,R_3,\dots)$.  It is atypical iff
$R_1\le m$; when typical, its dimension is $2^{m+1}\dim L_{\gl(m+1)}(\mu^T)$.  Hence:
\begin{itemize}
\item[(a)] is the atypical region, where $V_R$ already lives over $\gl(0|m)$, i.e.\ $\mathfrak{sl}_m$ duality;
\item[(b)] is $V_R=V_{R'}\otimes\mathrm{Ber}$, a twist by the one-dimensional Berezinian;
\item[(c)] for $R_2>m+1$, $V_R=0$ over $\gl(1|m+1)$, and the first algebra that sees $R$ is $\gl(2|m+2)$.
\end{itemize}
The Berezinian of $\gl(k|m+k)$ gives the rules
\[
  (b_k):\qquad R_{k+1}=m+k\ \Rightarrow\ H_R=H_{(R_1+1,\dots,R_k+1,R_{k+2},\dots)}\quad\text{on }A=q^{-m}.
\]
All $24$ instances with $|R|,|R'|\le6$ hold on every knot with data, checked modulo $2^{61}-1$
(\texttt{data/nonrect\_rule\_bk.tsv}).  The single instance with $k=2$ is $[2,2,2]\equiv[3,3]$ at $A=1$.

\subsection{Coloured Links--Gould polynomials and genus detection}
For $R(r)=(r,\mu)$ and $r>m$, $V_{R(r)}$ is the typical module of highest weight $(r\,|\,\mu^T)$.  For $m=1$ and
$\mu=(1^n)$ it is the $4(n+1)$-dimensional $\mathfrak{sl}(2|1)$-module $V(n,a)$ of \cite{LNG}, and $V(0,a)$ is
Links--Gould.  Define $P^\mu_m=H_{(r,\mu)}(q^{-m},q)|_{q^r=\lambda}$ and $e^\mu_m=\frac14\mathrm{span}_\lambda P^\mu_m$.
By \cite{LNG}, $(m+1)d\le e^\mu_m\le(m+1)g$.
\begin{observation}\label{obs:hooks}
\begin{itemize}
\item[(i)] $P^\mu_m$ is a Laurent polynomial in $(q,\lambda^2)$ with $P^\mu_m(1,\lambda)=\Delta(\lambda^2)^{m+1}$.
It is exact, by discrete Fourier reconstruction, in $12$ cases on $3_1,4_1,6_2,12n_{750}$ with $m\le2$.
It satisfies $P^\mu_m(q,\lambda)=P^{\mu^*}_m(q,q^{m-\ell(\mu)}/\lambda)$, where $\mu^*$ is the complement of $\mu$ in
the $\ell(\mu)\times(m+1)$ box.
\item[(ii)] $e^{[1]}_1=e^{[1,1]}_1=2g$ on all $60$ knots up to $10$ crossings of braid index $\le3$.
\item[(iii)] On the anomalous knots, see Table~\ref{tab:hooks}.  These values come from the vertex model
($r$ up to $22$--$58$) with a least-squares fit in $\lambda$, the fit window wider than the span and the support
confirmed with a second window; even powers only; agreement with the data where they exist $\le2\cdot10^{-10}$.
Braid index $5$--$6$ knots use an evaluator that works weight space by weight space.  For $\mu=[2]$ at $m=2$ a few
values of $r$ where the module construction is numerically inaccurate (errors up to $4\cdot10^{-5}$) were
discarded; with them included the fit produced spurious coefficients above the genus bound.
\end{itemize}
\end{observation}

\begin{table}[h]
\centering
\begin{tabular}{lcc|ccc|ccc}
\toprule
knot & $g$ & $d$ & $e_1$ & $e^{[1]}_1$ & $e^{[1,1]}_1$ & $e_2$ & $e^{[1]}_2$ & $e^{[2]}_2$\\
\midrule
$11n_{34}$ & 3 & 0 & 3 & 6 & 6 & 5 & 9 & 9\\
$11n_{42}$ & 2 & 0 & 3 & 4 & 4 & 5 & 6 & 6\\
$12n_{313},\ 12n_{430}$ & 2 & 0 & 3 & 4 & -- & 5 & 6 & --\\
$11n_{67},\ 11n_{97}$ (Type I) & 2 & 1 & 3 & 4 & -- & 4 & 6 & 6\\
$12n_{51},\ 12n_{268}$ (Type I) & 2 & 1 & 3 & 4 & -- & -- & -- & --\\
$11n_{73}$ & 3 & 2 & 4 & 6 & 6 & 6 & 9 & 9\\
$12n_{293},\ 12n_{321}$ (Type II) & 2 & 1 & 4 & 4 & 4 & 6 & 6 & 6\\
$12n_{457}$ (Type II) & 2 & 1 & 4 & 4 & -- & 6 & 6 & 6\\
$12n_{750},\ 12n_{830}$ & 3 & 2 & 6 & 6 & 6 & 9 & -- & --\\
\bottomrule
\end{tabular}
\caption{Symmetric (Links--Gould) versus hook towers at $m=1$ ($\gl(1|2)$) and $m=2$ ($\gl(1|3)$).  The genus bound is $(m+1)g$.}\label{tab:hooks}
\end{table}

\paragraph{Interpretation.}  The first hook families attain the genus bound $e^\mu_m=(m+1)g$ on every knot
tested, for $m=1$ and $m=2$, including all $14$ knots with $e_1<2g$ and every knot where the Links--Gould member
falls short at $m=2$.  So the anomalous defect is a property of the member
$\mu=\emptyset$, not of the knot.  This is the positive answer to \cite[Question~1.4]{LNG} on these examples.
Moreover $e^{[1]}_1$ distinguishes the mutants $11n_{34}$ and $11n_{42}$.  In the hook families the extra
$\lambda$-powers, whose coefficients vanish at $q=1$, appear already where the symmetric family has none.


\subsection{Two-row families ($k=2$)}
On $A=1$ ($\gl(2|2)$, $m=0$) the families $R=(r_2+\ell,r_2)$, $\ell=0,\dots,3$, with $r_2\ge2$ lie outside every
$\gl(1|1+\cdot)$ hook.  For $3_1,4_1,5_2$ the vertex model gives $\dim V_R=16(\ell+1)$ and agrees with the data on
$16$ checks ($\le4\cdot10^{-10}$).  At fixed $\ell$, $H_R$ is a Laurent polynomial in $\lambda=q^{r_2}$ (hold-out
error $\le5\cdot10^{-8}$ for $\ell\le2$), with $\frac14\mathrm{span}_\lambda=4g$ in all $12$ cases and the symmetry
$\lambda\mapsto q^{-\ell}/\lambda$, i.e.\ $(\lambda_1,\lambda_2)\mapsto(\lambda_2^{-1},\lambda_1^{-1})$.  For $3_1$ and
$R=(r,r)$ the polynomial is exact: a Laurent polynomial in $(q^2,\lambda^2)$ with value $\Delta(\lambda^2)^4$ at
$q=1$ (four odd roots of $\gl(2|2)$), matching $H_{[2,2]},H_{[3,3]}$.  Polynomiality jointly in $(\lambda_1,\lambda_2)$
could not be decided with $\ell\le3$.

\subsection{A definition of the differential expansion for arbitrary $R$}
On $A=q^{-m}$ let $k(R,m)=\min\{k\ge0:R_{k+1}\le m+k\}$.  Then $H_R$ is the invariant of the typical
$\gl(k|m+k)$-module $V_R$ with $k=k(R,m)$:
\begin{itemize}
\item $k=0$ is rule (a);
\item $k=1$ is the Links--Gould line;
\item $k\ge2$ is the region with no universal reduction.
\end{itemize}
We propose that the differential expansion of $R$ is the expansion of these multivariable invariants in
$\lambda_i=q^{R_i}$, $i\le k$.  It is organised by families $(r_1,\dots,r_k,\mu)$, each with its own tower
$e^\mu_m\in[(m+1)d,(m+1)g]$.  The universal coincidences are, conjecturally, exactly (a) and $(b_k)$.

%======================================================================
\section{Open questions}
\begin{enumerate}
\item Write out Steps 1--3 of Theorem~\ref{thm:lg} in full, and compare with an independent evaluation of
$LG^{2,1}$.
\item What selects the slope of the Links--Gould tower inside $[d,g]$?  Is $e^{[1]}_m=(m+1)g$ for all knots and
all $m$ (true for $m=1,2$ on every knot tested)?
\item Joint polynomiality in $(\lambda_1,\dots,\lambda_k)$ and towers for $k\ge2$, and $m\ge1$ for $k=2$.
\item Prove that (a) and $(b_k)$ are the only universal coincidences.
\item Prove Observation~\ref{obs:roots}.
\end{enumerate}

\begin{thebibliography}{99}
\bibitem{KM-defect} Ya.~Kononov, A.~Morozov, \emph{On the defect and stability of differential expansion},
arXiv:1504.07146.
\bibitem{KM-roots} Ya.~Kononov, A.~Morozov, \emph{Factorization of colored knot polynomials at roots of unity},
arXiv:1505.06170.
\bibitem{KT} B.-M.~Kohli, G.~Tahar, \emph{A lower bound for the genus of a knot using the Links--Gould invariant},
arXiv:2310.15617.
\bibitem{LNG} D.~L\'opez Neumann, S.~Garoufalidis, \emph{Genus bounds for knot polynomials of Lie superalgebras},
arXiv:2607.11735.
\bibitem{DKL} D.~De Wit, L.~H.~Kauffman, J.~R.~Links, \emph{The Links--Gould invariant of links},
arXiv:math/9811128.
\bibitem{MW} I.~Marin, E.~Wagner, \emph{A cubic defining algebra for the Links--Gould polynomial},
arXiv:1203.5981.
\end{thebibliography}

\end{document}
